\section{Experimental Design}
\label{sec:design}

\subsection{Track grid}

For each factor, the frozen agent signal $\SA$ (\S\ref{sec:pipeline}) is
executed under three configurations -- $\Cagent$, $\Ccz$, and $\Cstd$ --
producing $\Aref$, $\Acz$, and $\Ahxz$ respectively; alongside these,
$\CZ$ and $\HXZ$ are each team's own independently measured result, and
$\Pref$ is the paper's own reported number. Figure~\ref{fig:track-grid}
lays out all six reference points. Only the configuration varies across
$\Aref$, $\Acz$, and $\Ahxz$ -- the code hash of $\SA$ is identical across
all three by the freeze guarantee in \S\ref{sec:pipeline} -- so any
difference among them is attributable to configuration alone, never to a
silent change in the signal.

\begin{figure}[htbp]
  \centering
  \fbox{\parbox{0.9\linewidth}{\centering \vspace{4cm} F2: signal $\times$ config experiment grid, six reference points ($\Aref$/$\Acz$/$\Ahxz$/$\CZ$/$\HXZ$/$\Pref$) \vspace{4cm}}}
  \caption{The experiment grid. Only the configuration varies across tracks
    for a fixed factor; the frozen agent signal $S_A$ is held constant
    across $\Aref$, $\Acz$, and $\Ahxz$.}
  \label{fig:track-grid}
\end{figure}

\subsection{RQ1 design: the paper-anchored three-term identity}

$\Acz$ and $\Ahxz$ are constructed by taking the frozen $\SA$ and
re-executing it with $\Ccz$ or $\Cstd$ substituted for $\Cagent$ --
otherwise unchanged -- so that Equation~\eqref{eq:three-term}'s
configuration term, $\Acz - \Aref$ (or $\Ahxz - \Aref$), is measured
directly rather than inferred. The window/vintage residual term
(\S\ref{sec:framework}) is computed by requesting $\HXZ$ from
\texttt{hxz\_bridge.compute\_hxz\_reported\_both\_windows()} at the specific
sample window that matches $\Aref$/$\Ahxz$'s coverage, isolating how much of
the $\HXZ$-side gap would remain even if the two teams' data windows were
aligned.

\subsection{Field-level gap decomposition}
\begin{equation}
  \label{eq:gap-decomp}
  \sum \text{main effects} + \sum \text{interaction effects} = (\Acz - \Aref)
\end{equation}
Where the configuration fields that differ between $\Ccz$ and $\Cagent$
number five or fewer, a full factorial over those fields is run, and
Equation~\eqref{eq:gap-decomp} holds exactly -- every main effect and every
interaction term is directly measured, and the residual is exactly zero by
construction, not by assumption. Where the field count exceeds what a full
factorial can cover economically, a one-at-a-time (OAT) design is used
instead, at the cost of not measuring interaction terms explicitly.

\subsection{$t$-channel decomposition}
\begin{equation}
  \label{eq:t-channel}
  \log t = \log \mu - \log \sigma + \tfrac{1}{2} \log n
\end{equation}
This is an identity, not a model: for any track's reported mean return
$\mu$, return volatility $\sigma$, and sample size $n$, the three terms on
the right sum exactly to $\log t$. Unlike the mean return in
Equation~\eqref{eq:gap-decomp}, a $t$-statistic cannot be decomposed
additively across configuration fields, because it is a ratio rather than a
sum; Equation~\eqref{eq:t-channel} instead decomposes \emph{why} a $t$-statistic
changed between two tracks -- via a shift in mean return, in volatility, or
in sample size -- without requiring the mean-return decomposition's
factorial design.

\subsection{Identification levels and significance tiers}

Every pairwise comparison between two tracks is assigned an identification
level derived automatically from the resolved difference-set between their
two configurations, never hand-labeled: \emph{controlled}, when exactly one
configuration field differs; \emph{harmonized}, when the compared quantities
have been transformed onto a common contract before comparison;
\emph{observational}, when one side is an externally published number whose
underlying signal and data vintage cannot be controlled; and
\emph{unidentified}, when more than one field differs simultaneously.
Significance is reported at three tiers -- $|t| \geq 1.96$, $2.78$, and
$3.39$ -- matching the thresholds \citet{HouXueZhang2020} themselves use
under increasingly conservative multiple-testing corrections, rather than
a single conventional cutoff.

